\documentclass{cssheet}


%--------------------------------------------------------------------------------------------------------------
% Basic meta data
%--------------------------------------------------------------------------------------------------------------

\title{Endliche Gruppen}
\author{Prof. Dr. Christian Spannagel}
\date{\today}
\hypersetup{%
	pdfauthor={\theauthor},%
	pdftitle={\thetitle},% 
	pdfsubject={Aufgabenblatt Algebra},%
	pdfkeywords={algebra}
}

%--------------------------------------------------------------------------------------------------------------
% document
%--------------------------------------------------------------------------------------------------------------

\begin{document}
	\printtitle


\textbf{Aufgabe 1 (Zyklische Gruppen):} 

\begin{enumerate}[a)]
\item Analysiert die Gruppen, mit denen wir uns bisher beschäftigt haben. Welche davon sind zyklisch, welche nicht, und warum?
\item Beweist: Jede zyklische Gruppe ist abelsch.
\end{enumerate}

\textbf{Aufgabe 2 (Zyklische Gruppen und Isomorphie):} 

Versucht einmal, folgende Sätze zu beweisen! Tipp: Macht euch nochmal klar, was eine zyklische Gruppe ist, und sucht dann einen \emph{einfachen} Isomorphismus $\varphi$.
\begin{enumerate}[a)]
	\item Jede unendliche zyklische Gruppe ist isomorph zu $(\mathbb{Z};+)$.
	\item Jede endliche zyklische Gruppe der Ordnung $n$ ist isomorph zur Restklassengruppe $(\mathbb{Z}_n;\oplus)$.
\end{enumerate}

\textbf{Aufgabe 3 (Untergruppen):}  
\begin{enumerate}[a)]
\item Es gibt in jeder Gruppe mindestens zwei Untergruppen, die man \emph{triviale Untergruppen} nennt. Welche sind das? Beweist, dass es sich um Untergruppen handelt!
\item Beweist: Die Menge aller ganzzahligen Potenzen $a^g$ eines Elements einer Gruppe $G$ ist eine Untergruppe $U$ von $G$.
\item Analysiert die Gruppen, mit denen wir uns bisher beschäftigt haben, auf die Existenz nicht-trivialer Untergruppen. Welche könnt ihr finden? Welche Untergruppen sind davon zyklisch (mit welchen erzeugenden Elementen)?
\item Könnt ihr einen Zusammenhang entdecken zwischen der Ordnung einer Gruppe und der Ordnung einer Untergruppe dieser Gruppe? Was heißt dies für zyklische Gruppen?
\end{enumerate}


\vspace*{2cm}
\printlicense

\printsocials

\end{document}

\end{document}
